% Copia corregida de INFORME_CARRO_D_C_.pdf (primer corte).
% Cambios: motor unico (N20 6 V, 1:50, 500 RPM), masa medida de 317 g, radio de rueda
% unico y Ecs. 11 a 19 recalculadas. Los numeros del bloque VALUES salen de compute_values.py.
% Overleaf: subir este archivo y las 6 imagenes a la carpeta Imagenes/.
\documentclass[conference]{IEEEtran}
\usepackage[T1]{fontenc}
\usepackage[spanish,es-noshorthands,es-tabla]{babel}
\usepackage{graphicx}
\usepackage{amsmath}
\usepackage{booktabs}
\usepackage{siunitx}
\usepackage{url}
\sisetup{output-decimal-marker={,}, input-decimal-markers={.,}, group-separator={\,}, group-minimum-digits=5}
\DeclareSIUnit\kgf{kgf}
\DeclareSIUnit\rpm{RPM}
% BEGIN VALUES (generado por compute_values.py; no editar a mano)
\newcommand{\valMassKg}{0.317}
\newcommand{\valWheelRadiusM}{0.024}
\newcommand{\valWheelRadiusCm}{2.4}
\newcommand{\valRollingForceN}{0.0622}
\newcommand{\valRequiredTorqueNm}{7.463e-04}
\newcommand{\valRequiredTorqueKgfCm}{7.608e-03}
\newcommand{\valSafetyRated}{65.72}
\newcommand{\valSafetyStall}{262.9}
\newcommand{\valNoLoadRadS}{52.36}
\newcommand{\valMotorShaftRpm}{25000}
\newcommand{\valMotorShaftRadS}{2618}
\newcommand{\valSpeedNoLoad}{1.257}
\newcommand{\valSpeedRated}{1.005}
\newcommand{\valSpeedLimited}{1.047}
\newcommand{\valDutyFullCharge}{0.7143}
\newcommand{\valDutyAdopted}{0.6757}
\newcommand{\valEffectiveAtFull}{5.676}
\newcommand{\valOvervoltageNominalPct}{23.33}
\newcommand{\valOvervoltageFullPct}{40}
\newcommand{\valArmatureResistance}{10.91}
\newcommand{\valStallFullA}{0.77}
\newcommand{\valStallLimitedA}{0.4583}
\newcommand{\valStallTotalFullA}{1.54}
\newcommand{\valNoLoadUsePct}{3.5}
\newcommand{\valRatedUsePct}{7.5}
\newcommand{\valStallUsePct}{27.5}
\newcommand{\valStallFullUsePct}{38.5}
\newcommand{\valStallLimitedUsePct}{22.92}
\newcommand{\valBatteryNoLoadMeasured}{7.5}
\newcommand{\valBatteryLoadMeasured}{6.9}
\newcommand{\valBatteryNoLoadAbsError}{0.1}
\newcommand{\valBatteryNoLoadRelErrorPct}{1.351}
\newcommand{\valBatteryLoadAbsError}{0.5}
\newcommand{\valBatteryLoadRelErrorPct}{6.757}
% END VALUES

% Si la imagen no esta en el proyecto, deja un recuadro en lugar de detener la compilacion.
% Busca primero en Imagenes/ (Overleaf) y luego en figures/ (copia local).
% Argumento 2: nombre en Overleaf; argumento 3: nombre de la copia local.
\newcommand{\safeincludegraphics}[3][]{%
  \IfFileExists{Imagenes/#2}%
    {\includegraphics[#1]{Imagenes/#2}}%
    {\IfFileExists{figures/#3}%
      {\includegraphics[#1]{figures/#3}}%
      {\fbox{\parbox[c][2cm][c]{0.85\linewidth}{\centering\small [PENDIENTE: subir \detokenize{#2} a Imagenes/]}}}}}

% Marca visible para lo que falta por redactar o medir.
\newcommand{\pendiente}[1]{\textbf{[PENDIENTE: #1]}}

\begin{document}
\title{Vehículo DC con Motor N20 y Puente H L298N}
\author{Ricardo Cortes, Alejandro Moreno, Pablo Chavez y Michael~\pendiente{apellido}}
\maketitle

\begin{abstract}
Este documento presenta el diseño, el cálculo y la construcción de un vehículo de tracción eléctrica con dos motorreductores N20 de \SI{6}{V} y relación 1:50, controlados por un ESP32 a través de un puente H L298N y alimentados por dos celdas de litio en serie. Se calculan el torque requerido por la resistencia a la rodadura, la velocidad angular y lineal del vehículo y la corriente demandada en arranque y en régimen, contrastada con el límite del puente H. El torque nominal del motor supera unas \num[round-mode=places,round-precision=0]{\valSafetyRated} veces el requerido, la velocidad lineal teórica sin carga es de \SI{\valSpeedNoLoad}{m/s} a \SI{6}{V} y la corriente de arranque en el peor caso (\SI{\valStallFullA}{A}) queda por debajo de los \SI{2}{A} por canal del L298N. Para no superar el voltaje nominal del motor, el ciclo de trabajo del PWM se limita al \SI{67.5}{\percent}. Las mediciones sobre el vehículo construido muestran que el voltaje de las celdas cae de \SI{\valBatteryNoLoadMeasured}{V} a \SI{\valBatteryLoadMeasured}{V} con los motores en marcha. \pendiente{resultado de la velocidad medida y su error frente a la calculada}
\end{abstract}

\begin{IEEEkeywords}
Motor DC, puente H, L298N, N20, torque, relación de reducción, velocidad lineal, corriente de arranque, sobrevoltaje.
\end{IEEEkeywords}

\section{Introducción}
Los vehículos de tracción eléctrica representan un pilar fundamental en el aprendizaje de electrónica de potencia, control de motores y programación embebida. El proyecto consiste en construir un vehículo de corriente directa (DC) cuyo movimiento pueda ser dirigido por un microcontrolador ESP32 y, en una etapa posterior, integrado con el motor gráfico Unity mediante comunicación serial o inalámbrica (Wi-Fi/Bluetooth), aprovechando la conectividad nativa del ESP32.

El diseño aquí presentado emplea dos motorreductores N20 con reductor de engranajes metálico, un puente H dual L298N para el control de velocidad, y una fuente de alimentación compuesta por dos celdas de litio de \SI{3.7}{V} conectadas en serie. El vehículo completo, incluyendo chasis, ruedas, electrónica y baterías, presenta una masa total medida de $m = \SI{\valMassKg}{kg}$ (Fig.~\ref{fig:pesaje}).

El presente documento se estructura en seis secciones: introducción, marco teórico; luego se describe la metodología de diseño y los cálculos; posteriormente se analizan los resultados obtenidos (incluyendo la falla detectada); y finalmente se listan las referencias bibliográficas.

\section{Marco Teórico}
\subsection{Motor de Corriente Directa (DC)}
Un motor de corriente directa convierte energía eléctrica en energía mecánica mediante la interacción entre un campo magnético fijo (estátor) y una corriente que circula por el devanado del rotor. La relación fundamental entre el voltaje de alimentación $V$, la fuerza contraelectromotriz $\varepsilon$, la resistencia interna $R_a$ y la corriente $I$ es \cite{chapman}:
\begin{equation}
V = \varepsilon + I \cdot R_a
\end{equation}
En el momento del arranque, la fuerza contraelectromotriz es nula ($\varepsilon = 0$), por lo que la corriente alcanza su valor máximo (corriente de arranque o \emph{stall current}):
\begin{equation}
I_{\text{arranque}} = \frac{V}{R_a}
\end{equation}
El torque electromagnético producido por el motor es proporcional a la corriente \cite{chapman}:
\begin{equation}
\tau = K_\tau \cdot I
\end{equation}
donde $K_\tau$ es la constante de torque del motor, expresada en N$\cdot$m/A.

La fuerza contraelectromotriz es proporcional a la velocidad angular, $\varepsilon = K_e\,\omega$. Al combinar esta relación con las Ecs.~1 y~3 se obtiene la característica torque--velocidad del motor \cite{chapman}:
\[
\omega = \frac{V}{K_e} - \frac{R_a}{K_e K_\tau}\,\tau .
\]
La velocidad disminuye linealmente con la carga: es máxima sin carga ($\omega_0 = V/K_e$) y nula cuando el torque alcanza el valor de arranque ($\tau_{\text{arranque}} = K_\tau V/R_a$). Ambos extremos son proporcionales al voltaje aplicado, de modo que al reducir el voltaje efectivo mediante PWM se reducen en la misma proporción la velocidad sin carga y el torque disponible. Esta es la razón por la que la hoja de datos del motor reporta por separado los puntos sin carga, nominal y de arranque.

\subsection{Caja Reductora y Relación de Transmisión}
Una caja reductora de engranajes reduce la velocidad angular de salida en un factor $i$ (relación de reducción) y, en ausencia de pérdidas, amplifica el torque en el mismo factor \cite{shigley}. Considerando la eficiencia $\eta$ de la caja:
\begin{equation}
\omega_{\text{salida}} = \frac{\omega_{\text{motor}}}{i}
\end{equation}
\begin{equation}
\tau_{\text{salida}} = \tau_{\text{motor}} \cdot i \cdot \eta
\end{equation}

\subsection{Torque Requerido para el Desplazamiento}
Para que un vehículo se desplace sobre una superficie horizontal, el torque en las ruedas motrices debe vencer la resistencia al rodamiento. La fuerza de resistencia al rodamiento se calcula como \cite{gillespie}:
\begin{equation}
F_r = \mu_r \cdot m \cdot g
\end{equation}
donde $\mu_r$ es el coeficiente de resistencia al rodamiento (típicamente 0,01 a 0,02 para superficies duras), $m$ es la masa total del vehículo y $g = \SI{9.81}{m/s^2}$ es la aceleración gravitacional. El torque requerido en cada rueda motriz es:
\begin{equation}
\tau_{\text{req}} = F_r \cdot r_w
\end{equation}
siendo $r_w$ el radio de la rueda.

\subsection{Velocidad Lineal del Vehículo}
La velocidad lineal se obtiene a partir de la velocidad angular de la rueda:
\begin{equation}
v = \omega_{\text{salida}} \cdot r_w
\end{equation}
con $\omega_{\text{salida}}$ expresada en rad/s.

\subsection{Modulación por Ancho de Pulso (PWM) y Control de Voltaje}
Para controlar la velocidad de rotación de un motor de corriente directa sin necesidad de disipar energía en forma de calor a través de resistencias variables, se emplea la técnica de Modulación por Ancho de Pulso (\emph{Pulse Width Modulation}, PWM) \cite{chapman}. Esta técnica consiste en conmutar la tensión de alimentación a una frecuencia fija, variando el tiempo en que el pulso se encuentra en estado alto ($t_{\text{on}}$) en comparación con el periodo total de la señal ($T$).

El cociente entre el tiempo de encendido y el periodo total se define como el ciclo de trabajo (\emph{duty cycle}, $D$):
\begin{equation}
D = \frac{t_{\text{on}}}{T} \times \SI{100}{\percent}
\end{equation}
Como consecuencia de esta conmutación rápida, el motor, debido a su inercia mecánica y a la inductancia de los devanados, responde al valor medio del voltaje aplicado. El voltaje efectivo ($V_{\text{efectivo}}$) entregado al motor se expresa mediante la relación lineal:
\begin{equation}
V_{\text{efectivo}} = D \cdot V_{\text{cc}}
\end{equation}
donde $V_{\text{cc}}$ es el voltaje de la fuente de alimentación de potencia.

En el contexto del presente proyecto, las celdas de litio en serie suministran $V_{\text{cc}} = \SI{7.4}{V}$ (hasta \SI{8.4}{V} cargadas) y el motorreductor N20 opera nominalmente a $V_{\text{nominal}} = \SI{6}{V}$. Para no superar el voltaje nominal ni siquiera con las celdas completamente cargadas, el ciclo de trabajo debe cumplir $D \le V_{\text{nominal}}/V_{\text{cc,max}}$. Se adoptó un ciclo de trabajo máximo más conservador, equivalente a un voltaje efectivo de \SI{5}{V} con las celdas a \SI{7.4}{V}:
\begin{equation}
\begin{aligned}
D &\le \frac{6\,\text{V}}{8{,}4\,\text{V}} \approx \num{\valDutyFullCharge} \\
D_{\max} &= \frac{5\,\text{V}}{7{,}4\,\text{V}} \approx \num{\valDutyAdopted}\ (\SI{67.5}{\percent})
\end{aligned}
\end{equation}
Con ese valor, el voltaje efectivo con las celdas cargadas es $\num{\valDutyAdopted}\times\SI{8.4}{V} = \SI{\valEffectiveAtFull}{V}$, siempre por debajo del nominal. De este modo, limitar por software el ciclo de trabajo al \SI{67.5}{\percent} evita la necesidad de reguladores de voltaje físicos adicionales, protegiendo los devanados contra sobrecalentamiento destructivo.

\subsection{Puente H L298N}
El L298N es un circuito integrado de doble puente H de alta corriente. Permite controlar de forma independiente la dirección y velocidad de dos motores DC mediante señales TTL provenientes de un microcontrolador \cite{l298}. Sus parámetros eléctricos más relevantes son:
\begin{itemize}
\item Tensión de alimentación máxima: \SI{46}{V}
\item Corriente continua por canal: \SI{2}{A}
\item Corriente pico no repetitiva (\SI{100}{ms}): \SI{3}{A}
\item Corriente pico repetitiva (\SI{10}{ms}, ciclo 80/20): \SI{2.5}{A}
\item Disipación máxima a $T_{\text{case}} = \SI{75}{\celsius}$: \SI{25}{W}
\end{itemize}
Cada puente H está formado por cuatro transistores que conectan los terminales del motor a la alimentación o a tierra. Las entradas IN1 e IN2 (IN3 e IN4 en el segundo canal) fijan el sentido de la corriente, y la entrada de habilitación EN recibe la señal PWM que regula su valor medio. La Tabla~\ref{tab:puenteh} resume los estados de un canal según la hoja de datos \cite{l298}.

\begin{table}[!t]
\caption{Estados de un canal del puente H L298N \cite{l298}}
\label{tab:puenteh}
\centering
\begin{tabular}{@{}cccl@{}}
\toprule
EN & IN1 & IN2 & Estado del motor \\
\midrule
H (PWM) & H & L & Gira en un sentido \\
H (PWM) & L & H & Gira en el sentido contrario \\
H & IN1 = IN2 & & Frenado rápido \\
L & X & X & Giro libre \\
\bottomrule
\end{tabular}
\end{table}

Al estar construido con transistores bipolares, el L298N introduce una caída de saturación entre la alimentación y el motor, que aumenta con la corriente \cite{l298}. Por esta razón el voltaje que recibe el motor es menor que el producto $D \cdot V_{\text{cc}}$ de la Ec.~10.

La condición de seguridad para el sistema es que la corriente de arranque del motor no supere la corriente continua admisible del puente H por canal.

\section{Diseño y Metodología}
\subsection{Componentes del Sistema}
El vehículo fue construido con los siguientes componentes principales:
\begin{itemize}
\item $2\times$ motorreductores N20 metálicos de 3 etapas, voltaje nominal \SI{6}{V} DC, relación 1:50 y \SI{500}{RPM} sin carga.
\item $1\times$ puente H L298N (módulo de doble canal).
\item $2\times$ celdas de litio \SI{3.7}{V}/\SI{2000}{mAh} en serie (total \SI{7.4}{V}).
\item $1\times$ microcontrolador ESP32 (módulo genérico con chip USB-UART CP2102 y conector USB Type-C), como generador de señales PWM y lógica de control, con conectividad Wi-Fi/Bluetooth integrada para la futura comunicación con Unity. Sus pines de control operan a nivel lógico de \SI{3.3}{V}, compatible con los umbrales TTL de entrada del L298N (véase la Fig.~\ref{fig:pinout}).
\item Chasis, ruedas (radio $r_w = \SI{\valWheelRadiusCm}{cm} = \SI{\valWheelRadiusM}{m}$) y rueda loca trasera.
\end{itemize}

\subsection{Vista General del Vehículo Ensamblado}
Las Figs.~\ref{fig:superior} y~\ref{fig:inferior} muestran el vehículo completamente ensamblado desde su vista superior e inferior, respectivamente, evidenciando la disposición del chasis, las ruedas, los motorreductores N20 y la ubicación del puente H L298N y las celdas de alimentación.

\begin{figure}[!t]
\centering
\safeincludegraphics[width=0.8\columnwidth,height=0.3\textheight,keepaspectratio]{Auto_superior.jpeg}{fig1_top_view.png}
\caption{Vista superior del vehículo ensamblado.}
\label{fig:superior}
\end{figure}

\textbf{Nota sobre el voltaje de operación:} las celdas en serie entregan un voltaje nominal de \SI{7.4}{V} y hasta \SI{8.4}{V} cuando están completamente cargadas. El motorreductor N20 tiene un voltaje nominal de \SI{6}{V} DC. Si se aplicara la tensión de las celdas sin PWM, esto constituiría un sobrevoltaje del \SI{\valOvervoltageNominalPct}{\percent} (\SI{7.4}{V}) y de hasta el \SI{\valOvervoltageFullPct}{\percent} (\SI{8.4}{V}) respecto al valor nominal, lo cual es relevante para el análisis de la falla reportada.
% PENDIENTE: la falla no se describe en ninguna seccion; falta la seccion de Resultados.

\begin{figure}[!t]
\centering
\safeincludegraphics[width=\columnwidth,height=0.25\textheight,keepaspectratio]{Auto_inferior.jpeg}{fig2_bottom_view.png}
\caption{Vista inferior del vehículo ensamblado.}
\label{fig:inferior}
\end{figure}

\begin{figure}[!t]
\centering
\safeincludegraphics[width=\columnwidth,height=0.3\textheight,keepaspectratio]{ESP32_DataSheet.png}{fig3_pinout.png}
\caption{Diagrama de distribución de pines (\emph{pinout}) de la placa de desarrollo ESP32 con chip CP2102 (38 pines) utilizada para la generación de señales PWM \cite{techshop}.}
\label{fig:pinout}
\end{figure}

\subsection{Medición Experimental de la Masa del Vehículo}
Para validar el valor de masa utilizado en los cálculos de torque (Ec.~12), se realizó el pesaje directo del vehículo completamente ensamblado (chasis, ruedas, electrónica y baterías) sobre una balanza digital, tal como se muestra en la Fig.~\ref{fig:pesaje}. La balanza indicó \SI{317}{g}, valor que se usa como masa del vehículo en los cálculos. El radio de las ruedas motrices, medido con regla, es de \SI{\valWheelRadiusCm}{cm}.

\begin{figure}[!t]
\centering
\safeincludegraphics[width=\columnwidth,height=0.25\textheight,keepaspectratio]{Peso_auto_total.jpeg}{fig4_weighing.png}
\caption{Pesaje del vehículo ensamblado sobre balanza digital.}
\label{fig:pesaje}
\end{figure}

\subsection{Módulo Puente H L298N y Diagrama de Conexiones}
Para la etapa de potencia se implementó el módulo comercial basado en el circuito integrado L298N, cuya distribución física de pines, terminales de alimentación y esquema de aplicación se obtuvo de guías técnicas especializadas \cite{handson}. La disposición de estos terminales se ilustra en la Fig.~\ref{fig:l298n}.

\begin{figure}[!t]
\centering
\safeincludegraphics[width=\columnwidth]{Puente-H.jpg}{fig5_l298n_module.png}
\caption{Esquemático de pines y diagrama de conexiones del módulo controlador de motor Puente H L298N \cite{handson}.}
\label{fig:l298n}
\end{figure}

\subsection{Conexiones y Estrategia de Control}
La Tabla~\ref{tab:pines} resume las conexiones entre el ESP32 y el resto del sistema. El ESP32 genera las señales PWM con su periférico LEDC a \SI{1}{kHz} y 8 bits de resolución, y las aplica a las entradas de habilitación del L298N. Todos los componentes comparten la misma tierra, condición necesaria para que las señales lógicas de \SI{3.3}{V} del ESP32 tengan una referencia común con el puente H.

\begin{table}[!t]
\caption{Conexiones del ESP32}
\label{tab:pines}
\centering
\begin{tabular}{@{}lll@{}}
\toprule
GPIO & Destino & Función \\
\midrule
25 / 26 / 27 & ENA / IN1 / IN2 & Motor izquierdo \\
14 / 33 / 32 & ENB / IN3 / IN4 & Motor derecho \\
13 & Servo & Dirección (\SI{50}{Hz}) \\
GND & GND del L298N & Tierra común \\
\bottomrule
\end{tabular}
\end{table}

Además del límite del ciclo de trabajo de la Ec.~11, el firmware aplica dos medidas derivadas de los cálculos de corriente. Primero, una rampa que lleva el ciclo de trabajo de 0 a \SI{100}{\percent} del límite en \SI{250}{ms}, de forma que el motor gana velocidad, y por tanto fuerza contraelectromotriz, antes de recibir el voltaje completo; así la corriente inicial queda por debajo del valor de arranque de la Ec.~2. Segundo, un mecanismo de seguridad que detiene los motores si pasan \SI{300}{ms} sin recibir una orden válida, para que una pérdida de comunicación no deje el vehículo en movimiento.

\subsection{Parámetros del Motor (Datasheet)}
Los valores extraídos de la hoja de datos del fabricante \cite{foneacc} para el motorreductor N20 de \SI{6}{V} con caja reductora 1:50 y \SI{500}{RPM} sin carga se ilustran en la Fig.~\ref{fig:datasheet} y se resumen en la Tabla~\ref{tab:motor}. Las velocidades de la hoja de datos corresponden al eje de salida del reductor.

\begin{table}[!t]
\caption{Parámetros del motorreductor N20 (\SI{6}{V}, relación 1:50)}
\label{tab:motor}
\centering
\begin{tabular}{@{}lrl@{}}
\toprule
Parámetro & Valor & Unidad \\
\midrule
Voltaje nominal & 6 & V \\
Velocidad sin carga (salida) & 500 & RPM \\
Corriente sin carga & 0,07 & A \\
Velocidad nominal (salida) & 400 & RPM \\
Corriente nominal & 0,15 & A \\
Torque nominal & 0,5 & kgf$\cdot$cm \\
Corriente de arranque & 0,55 & A \\
Torque de arranque & 2 & kgf$\cdot$cm \\
Relación de reducción & 50 & --- \\
\bottomrule
\end{tabular}
\end{table}

\begin{figure*}[!t]
\centering
\safeincludegraphics[width=0.62\textwidth,height=0.45\textheight,keepaspectratio]{NOMBRE-MOTOR-POLOLU}{fig6_motor_datasheet.png}
\caption{Especificaciones técnicas, dimensiones y curvas de operación del motorreductor N20 extraídas de la hoja de datos del fabricante \cite{foneacc}. El motor del vehículo corresponde a la fila 1:50 con \SI{500}{RPM} sin carga.}
\label{fig:datasheet}
\end{figure*}

\subsection{Cálculo del Torque Requerido}
Con los datos del vehículo ($m = \SI{\valMassKg}{kg}$, $r_w = \SI{\valWheelRadiusM}{m}$) y un coeficiente de rodamiento $\mu_r = 0{,}02$ (superficie dura con posibles irregularidades):
\begin{equation}
F_r = 0{,}02 \times \SI{\valMassKg}{kg} \times \SI{9.81}{m/s^2} = \SI{\valRollingForceN}{N}
\end{equation}
El torque requerido por rueda motriz (2 ruedas en tracción diferencial, carga distribuida equitativamente):
\begin{equation}
\tau_{\text{req}} = \frac{F_r}{2}\cdot r_w = \frac{\SI{\valRollingForceN}{N}}{2}\times\SI{\valWheelRadiusM}{m} = \SI{\valRequiredTorqueNm}{N.m}
\end{equation}
Convirtiendo a kgf$\cdot$cm para comparar con la hoja de datos:
\begin{equation}
\begin{aligned}
\tau_{\text{req}} &= \SI{\valRequiredTorqueNm}{N.m}\times\frac{\SI{1}{kgf}}{\SI{9.81}{N}}\times\frac{\SI{100}{cm}}{\SI{1}{m}}\\
&\approx \SI{\valRequiredTorqueKgfCm}{\kgf.cm}
\end{aligned}
\end{equation}
El torque nominal del motorreductor a \SI{6}{V} es $\tau_{\text{salida}} = \SI{0.5}{\kgf.cm}$ (Tabla~\ref{tab:motor}), lo cual supera ampliamente el torque requerido:
\begin{equation}
\tau_{\text{salida}} = \SI{0.5}{\kgf.cm} \gg \tau_{\text{req}} = \SI{\valRequiredTorqueKgfCm}{\kgf.cm}
\end{equation}
El factor de seguridad en torque es:
\begin{equation}
FS = \frac{\tau_{\text{salida}}}{\tau_{\text{req}}} = \frac{0{,}5}{\num{\valRequiredTorqueKgfCm}} \approx \num{\valSafetyRated}
\end{equation}
Esto confirma que el motor seleccionado es adecuado para la carga propuesta, con un margen de seguridad muy superior a 4. Con el torque de arranque (\SI{2}{\kgf.cm}) el factor sería de aproximadamente \num{\valSafetyStall}. El cálculo solo considera la resistencia al rodamiento sobre superficie horizontal y a velocidad constante; no incluye la aceleración ni las pérdidas por fricción en el eje y en los engranajes.

\subsection{Cálculo de Velocidad Angular y Lineal}
La velocidad angular sin carga del eje de salida del reductor a \SI{6}{V} es de \SI{500}{RPM} (Tabla~\ref{tab:motor}). Convirtiendo a rad/s, y obteniendo la velocidad del eje del motor con la relación de reducción $i = 50$ (Ec.~4):
\begin{equation}
\begin{aligned}
\omega_{\text{salida}} &= 500\,\frac{\text{rev}}{\text{min}}\times\frac{2\pi\,\text{rad}}{1\,\text{rev}}\times\frac{1\,\text{min}}{60\,\text{s}} = \SI{\valNoLoadRadS}{rad/s}\\
\omega_{\text{motor}} &= i\cdot\omega_{\text{salida}} = \SI{\valMotorShaftRadS}{rad/s}\ (\SI{\valMotorShaftRpm}{RPM})
\end{aligned}
\end{equation}
La velocidad lineal teórica sin carga del vehículo resulta:
\begin{equation}
v = \omega_{\text{salida}} \times r_w = \SI{\valNoLoadRadS}{rad/s}\times\SI{\valWheelRadiusM}{m} = \SI{\valSpeedNoLoad}{m/s}
\end{equation}
Esa es la velocidad máxima teórica bajo condiciones sin carga a \SI{6}{V}. A la velocidad nominal bajo carga (\SI{400}{RPM}) la velocidad lineal es de \SI{\valSpeedRated}{m/s}. Con el voltaje efectivo limitado a \SI{5}{V} (Ec.~11), y suponiendo que la velocidad sin carga es aproximadamente proporcional al voltaje aplicado (Ec.~1), se espera una velocidad sin carga de \SI{\valSpeedLimited}{m/s}.

\subsection{Análisis de Corriente vs. Límite del Puente H}
La condición más exigente ocurre en el arranque, cuando ambos motores demandan simultáneamente su corriente máxima. Según la hoja de datos del fabricante \cite{foneacc} para el motor N20 de \SI{6}{V} y \SI{500}{RPM}, la corriente de arranque (\emph{stall current}) es $I_{\text{stall}} = \SI{0.55}{A}$ por motor, lo que corresponde a una resistencia de armadura $R_a = 6/0{,}55 = \SI{\valArmatureResistance}{\ohm}$ (Ec.~2). El L298N debe suministrar esta corriente por cada canal de forma independiente. Si la tensión de las celdas completamente cargadas (\SI{8.4}{V}) se aplicara sin limitar el PWM, la corriente de arranque sería $\SI{8.4}{V}/\SI{\valArmatureResistance}{\ohm}$:
\begin{equation}
\begin{aligned}
I_{\text{arranque}} &= \SI{0.55}{A} < I_{\text{L298N,cont}} = \SI{2}{A}\\
I_{\text{arranque,8,4 V}} &= \SI{\valStallFullA}{A} < I_{\text{L298N,cont}} = \SI{2}{A}
\end{aligned}
\end{equation}
La corriente de arranque se encuentra holgadamente dentro del límite continuo del puente H (\SI{2}{A}) incluso en el peor caso, y por debajo del límite pico repetitivo (\SI{2.5}{A}). Con ambos motores en arranque a \SI{8.4}{V} la batería entrega \SI{\valStallTotalFullA}{A}. En régimen nominal, la corriente por motor es de \SI{0.15}{A}, muy por debajo del límite. Los valores se resumen en la Tabla~\ref{tab:corrientes}.

\begin{table}[!t]
\caption{Comparación de corrientes del motor (\SI{6}{V}, \SI{500}{RPM}) vs. límite del L298N}
\label{tab:corrientes}
\centering
\begin{tabular}{@{}lrr@{}}
\toprule
Condición & Motor (A) & Uso de \SI{2}{A} (\%) \\
\midrule
Sin carga (\SI{6}{V}) & 0,070 & \num{\valNoLoadUsePct} \\
Régimen nominal (\SI{6}{V}) & 0,150 & \num{\valRatedUsePct} \\
Arranque (\SI{6}{V}) & 0,550 & \num{\valStallUsePct} \\
Arranque con PWM limitado (\SI{5}{V}) & \num[round-mode=places,round-precision=3]{\valStallLimitedA} & \num[round-mode=places,round-precision=1]{\valStallLimitedUsePct} \\
Arranque sin PWM (\SI{8.4}{V}) & \num[round-mode=places,round-precision=3]{\valStallFullA} & \num[round-mode=places,round-precision=1]{\valStallFullUsePct} \\
\bottomrule
\end{tabular}
\end{table}

\section{Resultados}
Esta sección compara las magnitudes calculadas en la Sección~III con las medidas sobre el vehículo construido. El error absoluto y el error relativo de cada magnitud se definen como:
\begin{equation}
e_{\text{abs}} = \left|x_{\text{medido}} - x_{\text{calculado}}\right|,\qquad
e_{\text{rel}} = \frac{e_{\text{abs}}}{x_{\text{calculado}}}\times\SI{100}{\percent}
\end{equation}

\subsection{Procedimiento de Medición}
Las mediciones se hicieron sobre el vehículo ensamblado, con el firmware de prueba de tracción, que mueve ambos motores hacia adelante con un ciclo de trabajo fijo y sin depender de la comunicación con Unity. Los voltajes se midieron con un multímetro digital en modo de voltaje DC, conectado en paralelo: entre los bornes \texttt{+12V} y \texttt{GND} del L298N para las celdas, y entre las salidas OUT1--OUT2 y OUT3--OUT4 para cada motor. La masa se obtuvo con una balanza digital y el radio de las ruedas con una regla. Como el PWM conmuta a \SI{1}{kHz}, el multímetro muestra el valor medio del voltaje, que es la magnitud que describe la Ec.~10. \pendiente{modelo y resolución del multímetro, superficie de la prueba de velocidad, distancia recorrida y número de repeticiones}

\subsection{Velocidad Lineal}
\pendiente{distancia recorrida, tiempos medidos, velocidad media y dispersión entre repeticiones}

\subsection{Voltaje de las Celdas}
El voltaje de las celdas en serie se midió con multímetro entre los bornes \texttt{+12V} y \texttt{GND} del L298N. Sin carga, con los motores detenidos, marcó \SI{\valBatteryNoLoadMeasured}{V}; con ambos motores girando, \SI{\valBatteryLoadMeasured}{V}. La caída de \SI{0.6}{V} bajo carga se atribuye a la resistencia interna de las celdas.
\pendiente{voltaje medido en los bornes de cada motor, indicando el ciclo de trabajo aplicado}

\subsection{Comparación entre lo Calculado y lo Medido}
La Tabla~\ref{tab:comparacion} reúne los valores calculados y los medidos.

\begin{table}[!t]
\caption{Comparación entre valores calculados y medidos}
\label{tab:comparacion}
\centering
\begin{tabular}{@{}lcccc@{}}
\toprule
Magnitud & Calculado & Medido & $e_{\text{abs}}$ & $e_{\text{rel}}$ (\%) \\
\midrule
Celdas sin carga (V) & 7,4 & \num{\valBatteryNoLoadMeasured} & \num{\valBatteryNoLoadAbsError} & \num[round-mode=places,round-precision=2]{\valBatteryNoLoadRelErrorPct} \\
Celdas con carga (V) & 7,4 & \num{\valBatteryLoadMeasured} & \num{\valBatteryLoadAbsError} & \num[round-mode=places,round-precision=2]{\valBatteryLoadRelErrorPct} \\
Velocidad lineal (m/s) & \num[round-mode=places,round-precision=3]{\valSpeedLimited} & --- & --- & --- \\
\bottomrule
\end{tabular}
\end{table}
% PENDIENTE: confirmar a que condicion corresponde cada valor calculado (PWM al 67,5 %, sin carga o con carga) antes de medir.

\subsection{Análisis de Discrepancias}
\textbf{Voltaje de las celdas.} Sin carga, las celdas marcaron \SI{\valBatteryNoLoadMeasured}{V}, un \num[round-mode=places,round-precision=2]{\valBatteryNoLoadRelErrorPct}\,\% por encima del valor nominal. La diferencia es esperable: \SI{3.7}{V} por celda es el voltaje nominal de la química de litio, y el voltaje real depende del estado de carga, entre unos \SI{3}{V} descargada y \SI{4.2}{V} cargada por celda. Con los motores en marcha el voltaje cayó a \SI{\valBatteryLoadMeasured}{V}, un \num[round-mode=places,round-precision=2]{\valBatteryLoadRelErrorPct}\,\% por debajo del nominal. El modelo de la Sección~III supone una fuente ideal de \SI{7.4}{V}, pero las celdas y el cableado tienen una resistencia interna que provoca una caída proporcional a la corriente demandada. En consecuencia, el voltaje efectivo de la Ec.~10 durante la marcha es menor que el calculado con el valor nominal.

\textbf{Voltaje en los motores.} A la caída de las celdas se suma la caída de saturación de los transistores del L298N, descrita en la Sección~II, que el modelo de la Ec.~10 no incluye. \pendiente{voltaje medido en cada motor, ciclo de trabajo aplicado y su comparación con $D \cdot V_{\text{cc}}$}

\textbf{Velocidad.} \pendiente{diferencia entre la velocidad medida y la calculada, explicada por el menor voltaje efectivo, la fricción en ejes y engranajes, la carga real sobre el motor y la tolerancia del radio de la rueda}

\subsection{Falla Detectada}
\pendiente{descripción de la falla mencionada en la Introducción y en la nota de la Sección~III-B: qué ocurrió, en qué condiciones y su relación con el sobrevoltaje}

\begin{thebibliography}{8}
\bibitem{chapman} S.~J. Chapman, \emph{Máquinas Eléctricas}, 5.ª ed. McGraw-Hill, 2012.
\bibitem{shigley} R.~G. Budynas y J.~K. Nisbett, \emph{Diseño en Ingeniería Mecánica de Shigley}, 9.ª ed. McGraw-Hill, 2014.
\bibitem{gillespie} T.~D. Gillespie, \emph{Fundamentals of Vehicle Dynamics}. Society of Automotive Engineers, 1992.
\bibitem{l298} STMicroelectronics, ``L298 Dual Full-Bridge Driver --- Datasheet,'' Rev.~9, 2000. [En línea]. Disponible en: https://www.st.com/resource/en/datasheet/l298.pdf
\bibitem{techshop} TechShopBD, ``ESP32 Development Board CP2102 Chip 38 Pin.'' [En línea]. Disponible en: https://techshopbd.com/product/esp32-development-board-cp2102-chip-38-pin (consultado: 22 de septiembre de 2026)
\bibitem{handson} Handson Technology, ``L298N Dual H-Bridge Motor Driver --- Specification and Datasheet,'' 2023. [En línea]. Disponible en: \url{https://www.handsontec.com/dataspecs/L298N%20Motor%20Driver.pdf}
\bibitem{foneacc} FONEACC Co., Ltd., ``FA-WG1218-N20 Micro Worm Gear Reducer DC Motor --- Datasheet,'' 2023. [En línea]. Disponible en: https://www.foneacc.com
\bibitem{farina} A.~L. Farina, ``Corriente continua,'' \emph{Ingeniería Eléctrica}, nov. 2020. [En línea]. Disponible en: \url{https://editores-srl.com.ar/sites/default/files/alberto_farina_20201130_corriente_continua_0.pdf}
% [8] Farina no se cita en el texto: PENDIENTE citarla o retirarla.
\end{thebibliography}
\end{document}
